%20 min preso!
\documentclass[xcolor=table]{beamer}
\usepackage{beamerthemesplit}
\usepackage{wrapfig}
\usetheme{SPbGU}
\usepackage{pdfpages}
\usepackage{amsmath}
\usepackage{cmap}
\usepackage[T2A]{fontenc}
\usepackage[utf8]{inputenc}
\usepackage[english]{babel}
\usepackage{indentfirst}
\usepackage{amsmath}
\usepackage{tikz}
\usepackage{multirow}
\usepackage[noend]{algpseudocode}
\usepackage{algorithm}
\usepackage{algorithmicx}
\usepackage{fancyvrb}
\usetikzlibrary{calc}
\usetikzlibrary{shapes,arrows}
\usetikzlibrary{arrows,automata}
\usetikzlibrary{positioning}
\usetikzlibrary{fit}
%\usepackage{caption}
\usepackage{subcaption}

\usepackage{kbordermatrix}
\renewcommand{\kbldelim}{(} 
\renewcommand{\kbrdelim}{)}

\newcommand\mca{\multicolumn{1}{c}{\cellcolor{red}\textbf{\{a\}}}}
\newcommand\mcb{\multicolumn{1}{c}{\cellcolor{red}\textbf{\{b\}}}}

\usepackage{tabularx}
\newcolumntype{Y}{>{\raggedleft\arraybackslash}X}

\renewcommand{\thealgorithm}{}

\newtheorem{mytheorem}{Theorem}
\renewcommand{\thealgorithm}{}

\newcommand{\tikzmark}[1]{\tikz[overlay,remember picture] \node (#1) {};}
\def\Put(#1,#2)#3{\leavevmode\makebox(0,0){\put(#1,#2){#3}}}

\newcommand{\ltz}{$< 1$}


\tikzset{
    state/.style={
           rectangle,
           rounded corners,
           draw=black, very thick,
           minimum height=2em,
           inner sep=2pt,
           text centered,
           },
}

\beamertemplatenavigationsymbolsempty

\title[All-Path CFPQ]{Context-Free Path Querying with All-Path Semantics by Matrix Multiplication}
%\subtitle[YaccConstructor]{Parsing techniques for graph analysis}
% То, что в квадратных скобках, отображается в левом нижнем углу.
\institute[JetBrains Research]{
JetBrains Research, Programming Languages and Tools Lab  \\
Saint Petersburg University
}

% То, что в квадратных скобках, отображается в левом нижнем углу.
\author[Rustam Azimov]{\textbf{Rustam Azimov}, Ilya Epelbaum, Semyon Grigorev }

\date{June 20, 2021}

\begin{document}
{
\begin{frame}[fragile]
  \begin{table}
  \centering
  \begin{tabularx}{\linewidth}{YcX}
    \includegraphics[height=1.5cm]{pictures/jetbrainsResearch.pdf} \hfill
    & \begin{minipage}[t]{0.3\textwidth}\center \vspace{-1cm}  GRADES-NDA 2021
      \end{minipage}
    & \hfill \includegraphics[height=1.5cm]{pictures/SPbGU_Logo.png}
  \end{tabularx}
  \end{table}
  \titlepage
\end{frame}
}

\begin{frame} \frametitle{Context-Free Path Querying}
  \begin{minipage}[m]{0.45\linewidth}
  \raisebox{-0.5\totalheight}{\includegraphics[width=\textwidth]{pictures/hierarchical.pdf}}
  \end{minipage}\hfill
  \begin{minipage}[m]{0.5\linewidth}
  Navigation through a graph
  \begin{itemize}
        \item Are nodes A and B on the same level of hierarchy?
        \item Is there a path of form $\textbf{Up}^n \, \textbf{Down}^n$?
        \item Find all paths of form $\textbf{Up}^n \, \textbf{Down}^n$ which start from the~node A
  \end{itemize}

  \end{minipage}

  \end{frame}

%  \begin{frame}[fragile] \frametitle{Applicatipons}
%    \begin{itemize}
%      \item Static code analysis
%      \item Graph database querying
%      \item RDF analysis
%    \end{itemize}
%  \end{frame}

  \begin{frame}[fragile]
    \frametitle{Context-Free Path Querying: Relational Query Semantics}
    \begin{itemize}
      \item $\mathbb{G} = (\Sigma, N, P)$ --- context-free grammar in normal form
      \begin{itemize}
        \item $A \rightarrow B C$, where $A, B, C \in N$
        \item $A \rightarrow x$, where $A \in N, x \in \Sigma \cup \{\varepsilon\}$
        \item $L(\mathbb{G},A) = \{ \omega \mid A \Rightarrow^* \omega \}$
      \end{itemize}
      \pause
      \item $G = (V,E,L)$ --- directed graph
        \begin{itemize}
          \item $v \xrightarrow{l} u \in E$
          \item $L \subseteq \Sigma$
        \end{itemize}
        \pause
      %\item $p = v_0 \xrightarrow{l_0} v_1 \xrightarrow{l_1} \cdots \xrightarrow{l_{n-2}} v_{n-1} \xrightarrow{l_{n-1}} v_n$ --- path in $G$
      \item $\omega(\pi) = \omega(v_0 \xrightarrow{l_0} v_1 \xrightarrow{l_1} \cdots \xrightarrow{l_{n-2}} v_{n-1} \xrightarrow{l_{n-1}} v_n) = l_0 l_1 \cdots l_{n-1}$
      \pause
      \item $R_A = \{ (n, m) \mid \exists n \pi m$, such that $\omega(\pi) \in L(\mathbb{G},A)\}$
    \end{itemize}
  \end{frame}

  \begin{frame}[fragile] \frametitle{Matrix-Based Algorithm: Relational Query Semantics}
    	\begin{algorithm}[H]
    		\begin{algorithmic}[1]
    			\caption{Context-free path querying algorithm}
    			\label{lst:algo1}
    			\Function{evalCFPQ}{$D=(V,E,L), G=(\Sigma,N,P)$}
    			\State{$n \gets$ |V|}
    			\State{$T \gets \{T^{A_i} \mid A_i \in N, T^{A_i}$ is a matrix $n \times n$, $T^{A_i}_{k,l} \gets$ \texttt{false}\} }
    			\ForAll{$(i,x,j) \in E$, $A_k \mid A_k \to x \in P$}
    			%\Comment{Matrices initialization}
    			%\For{$A_k \mid A_k \to x \in P$}
    			{$T^{A_k}_{i,j} \gets \texttt{true}$}
    			%\EndFor
    			\EndFor
    			\ForAll{$A_k \mid A_k \to \varepsilon \in P$}
    			\ForAll{$i \in \{0,\ldots ,n-1\}$}
    			{$T^{A_k}_{i,i} \gets \texttt{true}$}
    			\EndFor
    			\EndFor
    			
    			\While{any matrix in $T$ is changing}
    			%\Comment{Transitive c	losure calculation}
    			\For{$A_i \to A_j A_k \in P$}
    			{ $T^{A_i} \gets T^{A_i} + (T^{A_j} \times T^{A_k})$ } 
    			\EndFor
    			\EndWhile
    			\State \Return $T$
    			\EndFunction
    		\end{algorithmic}
    	\end{algorithm}
  \end{frame}


\begin{frame}[fragile]
\frametitle{Context-Free Path Querying: Single-Path Query Semantics}
\begin{itemize}
	\item $R_A = \{ (n, m) \mid \exists n \pi m$, such that $\omega(\pi) \in L(\mathbb{G},A)\}$ --- answers for the relational query semantics
	\pause
	\item For all $A \in N$, for all $(n,m) \in R_A$ also return some such path $n\pi m$
	\begin{itemize}
		\item usually the shortest path is returned
		\item returned path can be used as a proof of existence
	\end{itemize}
	\pause
	\item The main idea for the matrix-based algorithm is to store additional information in adjacency matrices to be able to restore one such path $n \pi m$ for all $(n,m) \in R_A$
	\begin{itemize}
		\item the intermediate vertex
		\item some additional information about path such as length
	\end{itemize}
	
\end{itemize}
\end{frame}

\begin{frame}[fragile]
	\frametitle{Context-Free Path Querying: All-Path Query Semantics}
	\begin{itemize}
		\item $R_A = \{ (n, m) \mid \exists n \pi m$, such that $\omega(\pi) \in L(\mathbb{G},A)\}$ --- answers for the relational query semantics
		\pause
		\item For all $A \in N$, for all $(n,m) \in R_A$ also return \textbf{all} such paths $n\pi m$
		\begin{itemize}
			\item the number of such paths can be infinite if the input graph has cycles
			\item usually the result is represented in the form of a finite structure such as annotated grammar or a parse forest
		\end{itemize}
		\pause
		\item Currently, our matrix-based algorithms cannot handle the all-path query semantics
		
		\item The only linear algebra-based algorithm that solves this problem is the Kronecker product-based CFPQ algorithm
		
	\end{itemize}
\end{frame}

\begin{frame}[fragile]
	\frametitle{Kronecker product-based algorithm}
	\begin{itemize}
		\item Does not require transformation of the input grammar
		\item Uses Recursive State Machines for the context-free grammar representation
		
	\end{itemize}

	\begin{figure}[h]
		\begin{tikzpicture}[shorten >=1pt,auto]
			\node[state, initial] (q_0)   {$q_S^0$};
			\node[state] (q_1) [right=of q_0] {$q_S^1$};
			\node[state] (q_2) [right=of q_1] {$q_S^2$};
			\node[state, accepting] (q_3) [right=of q_2] {$q_S^3$};
			\path[->]
			(q_0) edge node {a} (q_1)
			(q_1) edge node {S} (q_2)
			(q_2) edge node {b} (q_3)
			(q_1) edge [bend left = 38, above]  node {b} (q_3);
			\node[draw=black, fit= (q_0) (q_1) (q_2) (q_3), xshift=-4.5ex, inner sep=0.95cm, label=Box S] {};
		\end{tikzpicture}
		\centering
		\caption{The RSM for the grammar $S \to a S b \mid a b$ of the same-generation query}
	\end{figure}
\end{frame}

\begin{frame}[fragile]
	\frametitle{Kronecker product-based algorithm}
	\begin{itemize}
		\item Does not require transformation of the input grammar
		\item Uses Recursive State Machines for the context-free grammar representation
		
	\end{itemize}
	
	\begin{figure}[h]
		\begin{tikzpicture}[shorten >=1pt,auto]
			\node[state, initial] (q_0)   {$q_S^0$};
			\node[state] (q_1) [right=of q_0] {$q_S^1$};
			\node[state] (q_2) [right=of q_1] {$q_S^2$};
			\node[state, accepting] (q_3) [right=of q_2] {$q_S^3$};
			\path[->]
			(q_0) edge node {a} (q_1)
			(q_1) edge node {S} (q_2)
			(q_2) edge node {b} (q_3)
			(q_1) edge [bend left = 38, above]  node {b} (q_3);
			\node[draw=black, fit= (q_0) (q_1) (q_2) (q_3), xshift=-4.5ex, inner sep=0.95cm, label=Box S] {};
		\end{tikzpicture}
		\centering
		\caption{The RSM for the grammar $S \to a S b \mid a b$ of the same-generation query}
	\end{figure}
\end{frame}

\begin{frame}[fragile]
	\frametitle{Kronecker product-based algorithm}
	\begin{itemize}
		\item The algorithm computes the Kronecker product of the adjacency matrices of the graph for RSM and the input graph
		{\scriptsize
			$$
			\begin{pmatrix}
				. & \{a\} & . & .     \\
				. & . & \{S\} & \{b\} \\
				. & . & . & \{b\}     \\
				. & . & . & .
			\end{pmatrix}
			\otimes
			\begin{pmatrix}
				. & \{a\} & . & .     \\
				. & . & \{a\} & .     \\
				\{a\} & . & . & \{b\} \\
				. & . & \{b\} & .
			\end{pmatrix}
			=$$
		}
		{\scriptsize
			\renewcommand{\arraystretch}{0.5}
			\setlength\arraycolsep{0.1pt}
			\begin{align*}
				& \kbordermatrix{
					& (0,0) & (0,1) & (0,2) & (0,3) & \vrule & (1,0) & (1,1) & (1,2) & (1,3) & \vrule &  (2,0) & (2,1) & (2,2) & (2,3) & \vrule &  (3,0) & (3,1) & (3,2) & (3,3) &\\ 
					(0,0) & . & . & . & . & \vrule & . & \{a\} & . & . & \vrule & . & . & . & . &  \vrule & . & . & . & . \\
					(0,1) & . & . & . & . & \vrule & . & . & \mca & . & \vrule & . & . & . & . &  \vrule & . & . & . & . \\
					(0,2) & . & . & . & . & \vrule & \{a\} & . & . & . & \vrule & . & . & . & . &  \vrule & . & . & . & . \\
					(0,3) & . & . & . & . & \vrule & . & . & . & . & \vrule & . & . & . & . &  \vrule & . & . & . & . \\
					\hline
					(1,0) & . & . & . & .  & \vrule & . & . & . & . & \vrule & . & . & . & . & \vrule & . & . & . & . \\
					(1,1) & . & . & . & .  & \vrule & . & . & . & . & \vrule & . & . & . & . & \vrule & . & . & . & . \\
					(1,2) & . & . & . & .  & \vrule & . & . & . & . & \vrule & . & . & . & . & \vrule & . & . & . & \mcb \\
					(1,3) & . & . & . & .  & \vrule & . & . & . & . & \vrule & . & . & . & . & \vrule & . & . & \{b\} & . \\
					\hline
					(2,0) & . & . & . & .  & \vrule & . & . & . & . & \vrule & . & . & . & . & \vrule & . & . & . & . \\
					(2,1) & . & . & . & .  & \vrule & . & . & . & . & \vrule & . & . & . & . & \vrule & . & . & . & . \\
					(2,2) & . & . & . & .  & \vrule & . & . & . & . & \vrule & . & . & . & . & \vrule & . & . & . & \{b\} \\
					(2,3) & . & . & . & .  & \vrule & . & . & . & . & \vrule & . & . & . & . & \vrule & . & . & \{b\} & . \\
					\hline
					(2,0) & . & . & . & .  & \vrule & . & . & . & . & \vrule & . & . & . & . & \vrule & . & . & . & . \\
					(2,1) & . & . & . & .  & \vrule & . & . & . & . & \vrule & . & . & . & . & \vrule & . & . & . & . \\
					(2,2) & . & . & . & .  & \vrule & . & . & . & . & \vrule & . & . & . & . & \vrule & . & . & . & . \\
					(2,3) & . & . & . & .  & \vrule & . & . & . & . & \vrule & . & . & . & . & \vrule & . & . & . & . \\
				}
			\end{align*}
		}
		\end{itemize}
\end{frame}

\begin{frame}[fragile] \frametitle{Research Questions}
\begin{itemize}
	\item Can we extend the matrix-based CFPQ algorithm to all-path query semantics?
	\item What the cost of such extension?
	\item How does the matrix-based solution for the all-path query semantics compare to the Kronecker product-based?
\end{itemize}
\end{frame}

\begin{frame}[fragile] \frametitle{All-Path Index}
\begin{itemize}
	\item We store the additional information about the paths found as the sets of the intermediate vertices
	\item We introduce the following matrix multiplication operation
	\item $T^A \odot T^B = T^C$ where $T^C_{i,j} = \bigcup^{n}_{k = 1} (T^A_{i,k} \otimes T^B_{k,j})$ and  
\end{itemize}
$$T^A_{i,k} \otimes T^B_{k,j} = \begin{cases} \{k\}, & \mbox{if } T^A_{i,k} \neq \emptyset \wedge T^B_{k,j} \neq \emptyset \\ \emptyset, & \mbox{otherwise} \end{cases}$$
\end{frame}

  \begin{frame}[fragile] \frametitle{Matrix-Based Algorithm: All-Path Query Semantics}
\begin{algorithm}[H]
	\begin{algorithmic}[1]
		\caption{CFPQ algorithm w.r.t. all-path query semantics}
		\label{lst:algo2}
		\Function{AllPathCFPQ}{\par
			\hskip\algorithmicindent $D = (V, E, \Sigma)$, \par
			\hskip\algorithmicindent $G=(N,\Sigma,P,S)$} %\Comment{Grammar in WCNF}\par
		\State{$n \gets$ |V|}
		\State{$T \gets \{T^{A} \mid A \in N, T^{A}$ is a matrix $n \times n$, $T^{A}_{i,j} \gets \emptyset$ \} }
		\ForAll{$(i,x,j) \in E$, $A \mid A \to x \in P$}
		%\Comment{Matrices initialization}
		%\For{$A_k \mid A_k \to x \in P$}
		{$T^{A}_{i,j} \gets \{n\}$}
		%\EndFor
		\EndFor
		\ForAll{$A \mid A \to \varepsilon \in P$}
		{$T^{A}_{i,i} \gets \{n\}$}
		\EndFor
		
		\While{any matrix in $T$ is changing}
		%\Comment{Transitive closure calculation}
		\ForAll{$A \to B C \in P$ where $T^{B}$ or $T^{C}$ are changed}
		\State{ $T^{A} \gets T^{A} + (T^{B} \odot T^{C})$ } 
		\EndFor
		\EndWhile
		\State \Return $T$
		\EndFunction
	\end{algorithmic}
\end{algorithm}
\end{frame}

 % \begin{frame}[fragile] \frametitle{Matrix-Based Algorithm: Technical Details}
  %  \begin{itemize}
   %   \item We can remove $\textit{length}$ or $\textit{height}$ to reduce memory consumption
    %  \item The PathIndex operations can be represented as bitwise atomic operations
     % \item We still can use existing high-performance libraries for matrix operations if they support the creation of custom operations
    %\end{itemize}
  %\end{frame}


%\begin{frame}[fragile] \frametitle{Example: Graph and Grammar}
%  	\[
%  	\includegraphics[width=5cm]{pictures/example_graph.pdf}
%  	\]
%  	\pause
%  	\[
%  	\begin{array}{rcclcrccl}
%  	0: & S & \rightarrow & A \ B   & \quad & 3: & A & \rightarrow & \text{\emph{a}}     \\
%  	1: & S & \rightarrow & A \ S_1       & \quad & 4: & B & \rightarrow & \text{\emph{b}} \\
%  	2: & S_1 & \rightarrow & S \ B & & & & &
%  	
%  	\end{array}
%  	\]
%  	
%\end{frame}

%\begin{frame}[fragile] \frametitle{Example: Initial Matrices}
%\[
%T^{(1),A} = \begin{pmatrix}
%\bot & (0,1,0,1,1)       & \bot & \bot       \\
%\bot & \bot & (1,2,1,1,1)       & \bot \\
%(2,0,2,1,1)       & \bot & \bot & \bot \\
%\bot       & \bot & \bot & \bot \\
%\end{pmatrix}
%\]
%\[
%T^{(1),B} = \begin{pmatrix}
%\bot & \bot       & \bot & (0,3,0,1,1)       \\
%\bot & \bot & \bot       & \bot \\
%\bot       & \bot & \bot & \bot \\
%(3,0,3,1,1)      & \bot & \bot & \bot \\
%\end{pmatrix}
%\]
%
%\end{frame}

%\begin{frame}[fragile] \frametitle{Example: Final Matrix}
%\[
%T^{(14),S} = \begin{pmatrix}
%(0,0,1,12,12) & \bot       & \bot & (0,3,1,6,6)       \\
%(1,0,2,4,4) & \bot & \bot       & (1,3,2,10,10) \\
%(2,0,0,8,8)       & \bot & \bot & (2,3,0,2,2) \\
%\bot       & \bot & \bot & \bot \\
%\end{pmatrix}
%\]
%
%\end{frame}

%\begin{frame}[fragile] \frametitle{Path Extraction Algorithm}
%\begin{algorithm}[H]
%	\begin{algorithmic}[1]
%		\caption{Path extraction algorithm}
%		\label{lst:algo3}
%		\Function{extractPath}{$i, j, A, T=\{T^{A_i}\}, G=(N,\Sigma,P)$}
%		\State{$index \gets T^{A}_{i,j}$ }
%		
%		\If{$index = \bot$}
%		\State \Return $\pi_{\emptyset}$
%		\Comment{Such a path does not exist}
%		\EndIf
%		
%		\If{$index.height = 1$}
%		\If{$index.length = 0$}
%		\State \Return $[]$
%		\Comment{Return an empty path}
%		\EndIf
%		\ForAll{$ x \mid (i,x,j) \in E$}
%		\If{$A \to x \in P$}
%		\State \Return $[(i,x,j)]$
%		\Comment{Return a path of length one}
%		\EndIf
%		\EndFor
%		\EndIf
%		
%		\ForAll{$A \to B C \in P$}
%		\State{$index_B \gets T^{B}_{i,index.middle}$ }
%		\State{$index_C \gets T^{C}_{index.middle,j}$ }			
%		\If{$(index_B \neq \bot) \wedge (index_C \neq \bot)$}
%		\State{$maxH \gets max(index_B.height, index_C.height)$ }
%		\If{$index.height = maxH + 1$}
%		
%		
%		\State{$\pi_1 \gets$ \Call{extractPath}{$i, index.middle, B, T, G$}}
%		\State{$\pi_2 \gets$ \Call{extractPath}{$index.middle, j, C, T, G$}}
%		\State \Return $\pi_1 + \pi_2$
%		\Comment{Return the concatenation of paths}
%		\EndIf
%		\EndIf
%		\EndFor
%		\EndFunction
%	\end{algorithmic}
%\end{algorithm}
%\end{frame}


\begin{frame}[fragile] \frametitle{Path extraction}

\begin{itemize}
	\item After constructing a set of matrices with sets of intermediate vertices, we can extract all required paths $i\pi j$ for every vertex pair $i, j$ if such paths exist
	\item It is	assumed that the sets of paths are computed lazily, to ensure the	termination in case of an infinite number of paths
	
\end{itemize}
\end{frame}


\begin{frame}[fragile] \frametitle{Implementation}

\begin{itemize}
	\item For evaluation we use the following CPU-based implementations of CFPQ algorithms with sparse matrix representation
	\begin{itemize}
		\item \textbf{$MtxRel$} --- for  relational query semantics that uses \textbf{pygraphblas} --- a Python wrapper around the GraphBLAS API
		\item \textbf{$MtxSingle$} --- for  single-path query semantics that also uses \textbf{pygraphblas}
		\item \textbf{$MtxAll$} --- the implementation of the proposed matrix-based algorithm for all-path query semantics which utilizes \textbf{SuiteSparse} and our own Python wrapper
		\item \textbf{$Tns$} --- our Python implementation of the Kronecker product-based algorithm for all-path query semantics
		
	\end{itemize}
\end{itemize}
\end{frame}


\begin{frame} \frametitle{Evaluation}
  \begin{itemize}
   \item OS: Ubuntu 18.04
   \item CPU: Intel core i7 6700 3,4GHz
   \item RAM: DDR4 64 Gb
  \end{itemize}
\end{frame}

\begin{frame}[fragile] \frametitle{Evaluation: CFPQ\footnote{Time in seconds and memory is measured in megabytes}}
\begin{center}
	\tikzmark{yyy}{
	}
	{\scriptsize
	{\setlength{\tabcolsep}{0.3em}
			\rowcolors{3}{}{lightgray}
					\begin{tabular}{|l|l|l|l|l|l|l|l|l|l|l|}
						\hline
						\multicolumn{1}{|c|}{\multirow{2}{*}{Graph}} & \multicolumn{1}{c|}{\multirow{2}{*}{\#V}} & \multicolumn{1}{c|}{\multirow{2}{*}{\#E}} &  \multicolumn{2}{c|}{MtxRel} & \multicolumn{2}{c|}{MtxSingle} & \multicolumn{2}{c|}{MtxAll} & \multicolumn{2}{c|}{Tns} \\ \cline{4-11} 
						\multicolumn{1}{|c|}{}                       & \multicolumn{1}{c|}{}                     & \multicolumn{1}{c|}{}                     & Time         & Mem          & Time        & Mem        & Time           & Mem           & Time         & Mem    \\ \hline
						pathways                                     & 6 238                                     & 18 598 & 0.01         & 140  & 0.01           & 671 & 0.01         & 49           & 0.01        & 122               \\ \hline
						go-hierarchy                                 & 45 007                                    & 980 218                                   & 0.09         & 255 & 0.84           & 671 & 0.35        & 195        & 0.24        & 252                             \\ \hline
						enzyme                                       & 48 815                                    & 109 695                                   & 0.01         & 181 & 0.01           & 217 & 0.02         & 61           & 0.02        & 132                            \\ \hline
						eclass\_514en                                & 239 111                                   & 523 727                                   & 0.06         & 181 & 0.16           & 216    & 0.22         & 126          & 0.27        & 193                         \\ \hline
						go                                           & 272 770                                   & 534 311                                   & 0.94         & 246  & 0.93           & 217  & 1.13         & 990          & 1.27        & 243                          \\ \hline
						%geospecies                                   & 450 609                                   & 2 311 461                   & 0.01         & 248 & 0.01           & 2251            & 0.34         & 156          & 0.01        & 196              \\ \hline
						geospecies & 450 609                                   & 2 311 461 & 7.48         & 7645     & 15.54          & 22941 & 32.06        & 44235        & 26.32       & 19537 \\ \hline
						taxonomy                                     & 5 728 398                                 & 14 922 125                            & 0.72         & 1175      & 1.15           & 2250    & 3.84        & 1507        & 3.56        & 1776                   \\ \hline
					\end{tabular}
	}
}
 \end{center}
%\pause
%\onslide<2>{\tikz[overlay,remember picture]{\draw[draw=red,thick,fill opacity=0.2] %($ (yyy) + (3,0.85)$) rectangle ($ (yyy) + (7.18,-1.6)$);}}
%\pause
%\onslide<3>{\tikz[overlay,remember picture]{\draw[draw=red,thick,fill opacity=0.2] %($ (yyy) + (4.6,1.2)$) rectangle ($ (yyy) + (5.3,-1.58)$);}}
%\pause
%\onslide<4>{\tikz[overlay,remember picture]{\draw[draw=red,thick,fill opacity=0.2] %($ (yyy) + (5.4,1.2)$) rectangle ($ (yyy) + (6.1,-1.58)$);}}
%\pause
%\onslide<5>{\tikz[overlay,remember picture]{\draw[draw=red,thick,fill opacity=0.2] ($ (yyy) + (9.2,1.2)$) rectangle ($ (yyy) + (9.9,-1.58)$);}}
\end{frame}

\begin{frame}[fragile] \frametitle{Evaluation: Average Path Extraction Time For \textit{go}}
\begin{center}
    \[
		\includegraphics[width=10cm]{pictures/go_10_all_matrixall_tensor.pdf}
	\]
\end{center}

\end{frame}

\begin{frame}[fragile] \frametitle{Conclusion}
  \begin{itemize}
    \item We propose a matrix-based CFPQ algorithm for all-path query semantics
    \pause
    \item The proposed algorithm constructs index up to 2-3 times slower and consumes more memory than the algorithm for single-path query semantics
    \pause
    \item If it is necessary to frequently recalculate
    the index for a changing graph or a path query then the best
    choice is the Kronecker product-based algorithm with
    faster and less memory consuming index construction
    \pause
    \item If it is     necessary to extract paths many times for a once constructed
    index or index changes can be efficiently computed dynamically then the proposed matrix-based CFPQ algorithm is  preferable
  \end{itemize}
  %\pause
  %\begin{itemize}
   % \item Dataset is published: both graphs and queries
    %\begin{itemize}
    %	\item Link: \url{https://github.com/JetBrains-Research/CFPQ_Data}
    %\end{itemize}
    
    %\item Implementations are available on GitHub
    %\begin{itemize}
    %	\item Link: \url{https://github.com/JetBrains-Research/CFPQ_PyAlgo}
    %\end{itemize}
    
  %\end{itemize}
\end{frame}

\begin{frame}[fragile] \frametitle{Future Research}
  \begin{itemize}
  	\item We compare the CPU-based implementation. In the future, we
  	want to obtain GPU-based and distributed implementations 
  	\item Also, further improvements in index creation and path extraction
  	for both matrix-based and Kronecker product-based algorithms are
  	required
    \item We plan to provide the multiple-source
    modifications for all linear algebra-based CFPQ algorithms
\end{itemize}
\end{frame}

\begin{frame}
\frametitle{Contact Information}
\begin{itemize}
  \item Semyon Grigorev:
    \begin{itemize}
      \item \href{mailto:s.v.grigoriev@spbu.ru}{s.v.grigoriev@spbu.ru}
      \item \href{mailto:Semen.Grigorev@jetbrains.com}{Semen.Grigorev@jetbrains.com}
    \end{itemize}
  \item Rustam Azimov:
  \begin{itemize}
  	\item \href{mailto:rustam.azimov19021995@gmail.com}{rustam.azimov19021995@gmail.com}
  	\item \href{mailto:Rustam.Azimov@jetbrains.com}{Rustam.Azimov@jetbrains.com}
  \end{itemize}
  \item Ilya Epelbaum: \href{mailto:iliyepelbaun@gmail.com}{iliyepelbaun@gmail.com}
\vspace{0.5cm}
  \item Dataset: \href{https://github.com/JetBrains-Research/CFPQ_Data}{https://github.com/JetBrains-Research/CFPQ\_Data}
   \item Algorithm implementations: \href{https://github.com/JetBrains-Research/CFPQ_PyAlgo}{https://github.com/JetBrains-Research/CFPQ\_PyAlgo}
\end{itemize}
\vspace{0.1cm}
\center{\huge{Thanks!}}
\end{frame}
\end{document}
